\subsection{Preuve de la loi de réciprocité quadratique utilisant les résultants}

Rappelons que la loi de réciprocité quadratique est une égalité entre deux
signes ; l'idée est alors d'utiliser la relation suivante sur les résultants
\[
\operatorname{Res}(P,Q)
=
(-1)^{\deg P \times \deg Q}\operatorname{Res}(Q,P),
\]
pour des polynômes \(P\) et \(Q\) de degrés respectifs \((p-1)/2\) et
\((q-1)/2\), où \(p\) et \(q\) sont des premiers impairs distincts, de sorte que
\[
\operatorname{Res}(P,Q)=\left(\frac{p}{q}\right)
\qquad\text{et}\qquad
\operatorname{Res}(Q,P)=\left(\frac{q}{p}\right).
\]

\subsubsection{Lemme}

\textit{Pour tout \(p\) premier impair, il existe \(Q_p\in\mathbb{Z}[X]\) tel que}
\[
X^{p-1}+X^{p-2}+\cdots+X+1
=
X^{(p-1)/2}Q_p\left(X+\frac{1}{X}\right).
\]
\textit{En outre, modulo \(p\), on a}
\[
Q_p(Y)\equiv (Y-2)^{\frac{p-1}{2}}\pmod p.
\]

\begin{proof}
En posant \(Y=X^{-1}\), on cherche \(Q_p\) tel que
\[
Q_p(X+Y)
=
X^{(p-1)/2}+\cdots+X+1+Y+\cdots+Y^{(p-1)/2}.
\]
Le polynôme étant symétrique, il existe \(R\in\mathbb{Z}[X,Y]\) tel que le
membre de droite de l'égalité ci-dessus est égal à \(R(X+Y,XY)\). En notant
que \(XY=1\), il suffit alors de poser
\[
Q_p(X)=R(X,1).
\]
Raisonnons alors modulo \(p\) :
\begin{align*}
X^{p-1}+\cdots+X+1
&\equiv (X-1)^{p-1} \\
&\equiv (X^2-2X+1)^{(p-1)/2} \\
&\equiv
X^{(p-1)/2}
\left(X+\frac{1}{X}-2\right)^{(p-1)/2}
\pmod p.
\end{align*}
Soit
\[
Q_p\left(X+\frac{1}{X}\right)
\equiv
\left(X+\frac{1}{X}-2\right)^{(p-1)/2}
\pmod p,
\]
et
\[
Q_p(X)\equiv (X-2)^{(p-1)/2}\pmod p.
\]
\end{proof}

\subsubsection{Lemme}

\textit{Pour \(p\neq q\) des nombres premiers impairs, le résultant de
\(Q_p\) et \(Q_q\) est égal à \(\pm1\).}

\begin{proof}
Raisonnons par l'absurde et soit \(\ell\) premier divisant
\(\operatorname{Res}(Q_p,Q_q)\), de sorte que modulo \(\ell\), \(Q_p\) et
\(Q_q\) ont une racine commune
\[
\beta\in\mathbb{F}_{\ell^n},
\]
où l'entier \(n\) est inférieur ou égal à
\[
\min\left\{\frac{p-1}{2},\frac{q-1}{2}\right\}.
\]

Soit alors \(x \in \overline{\mathbb{F}}_\ell\) tel que
\[
x^2-\beta x+1=0.
\]
Alors
\[
x^{p-1}+\cdots+x+1
=
x^{(p-1)/2}\overline{Q}_p(\beta)
=
0.
\]

En multipliant cette égalité par \(x-1\), on en déduit que
\[
x^p=1
\]
dans \(\overline{\mathbb{F}}_\ell\). De la même façon, on a aussi
\[
x^q=1.
\]
Comme \(p\wedge q=1\), il vient \(x=1\), et donc
\[
p\equiv q\equiv 0 \pmod \ell,
\]
ce qui n'est pas possible, car \(p\wedge q=1\).
\end{proof}
